% ============================================================
%  Artículo Técnico — Análisis e Interpretación de Datos
%  UNIR — Actividad 2
%  Motor: MikTeX — compilar con pdflatex (2 pasadas)
% ============================================================
\documentclass[11pt, a4paper]{article}

% --- Codificación y lengua ---
\usepackage[utf8]{inputenc}
\usepackage[T1]{fontenc}
\usepackage[spanish]{babel}

% --- Fuente Georgia (aproximación con mathptmx + ajuste) ---
% MikTeX no incluye Georgia nativa; se usa TeX Gyre Termes,
% equivalente métricamente a Times New Roman / Georgia.
\usepackage{tgtermes}

% --- Márgenes y geometría ---
\usepackage[
  a4paper,
  top=2.5cm, bottom=2.5cm,
  left=2.5cm, right=2.5cm
]{geometry}

% --- Interlineado 1.5 ---
\usepackage{setspace}
\onehalfspacing

% --- Matemáticas y tablas ---
\usepackage{amsmath}
\usepackage{booktabs}
\usepackage{array}
\usepackage{float}

% --- Gráficos ---
\usepackage{pgfplots}
\pgfplotsset{compat=1.18}
\usepackage{tikz}

% --- Colores ---
\usepackage{xcolor}
\definecolor{azulunir}{RGB}{0, 84, 166}

% --- Hiperenlaces ---
\usepackage[
  colorlinks=true,
  linkcolor=azulunir,
  citecolor=azulunir,
  urlcolor=azulunir
]{hyperref}

% --- Bibliografía (estilo numérico estándar, sin natbib) ---

% --- Cabecera y pie de página ---
\usepackage{fancyhdr}
\pagestyle{fancy}
\fancyhf{}
\renewcommand{\headrulewidth}{0.4pt}
\fancyhead[L]{\small\textit{Análisis e Interpretación de Datos — UNIR}}
\fancyhead[R]{\small\thepage}

% --- Título compacto ---
\usepackage{titling}
\setlength{\droptitle}{-2em}

% ============================================================
\begin{document}

% ---- PORTADA / CABECERA DEL ARTÍCULO ----------------------
\begin{center}
  {\LARGE\bfseries Impacto del calor extremo sobre la mortalidad en España:\\
  Un análisis estadístico de las temporadas 2024--2025}

  \vspace{0.8em}
  {\large Autor: David Vicente Moya}\\[0.3em]
  {\normalsize Afiliación: Universidad Internacional de La Rioja (UNIR)}\\[0.3em]
  {\normalsize \today}
\end{center}

\vspace{0.5em}
\hrule
\vspace{0.8em}

% ---- RESUMEN ----------------------------------------------
\noindent\textbf{Resumen.}\quad
El presente artículo analiza el impacto del calor extremo sobre la mortalidad
en España durante las temporadas estivales de 2024 a 2025, empleando datos
semanales del Sistema de Monitorización de la Mortalidad Diaria (MoMo) del
Instituto de Salud Carlos III. Se aplican dos modelos estadísticos
complementarios: un análisis descriptivo de series temporales que caracteriza
la evolución del exceso de mortalidad a lo largo de las 68 semanas del periodo
de estudio, y una regresión lineal simple que cuantifica la relación entre las
defunciones atribuibles a la temperatura y el exceso de mortalidad observado
en las semanas estivales. Los resultados muestran un incremento sustancial de
la mortalidad atribuible al calor en el verano de 2025 respecto al de 2024
(3.660 frente a 2.008 defunciones), con un coeficiente de determinación
$R^2 = 0{,}57$ en la regresión lineal, lo que revela que la temperatura
explica una parte significativa, aunque no total, del exceso de mortalidad
observado. Se discuten las implicaciones para la salud pública y las
limitaciones de ambos enfoques.

\vspace{0.8em}
\hrule
\vspace{1em}

% ============================================================
\section{Introducción y estado del tema}
% ============================================================

El cambio climático ha intensificado la frecuencia e intensidad de las olas de
calor en la cuenca mediterránea, convirtiendo el exceso de temperatura en uno
de los principales riesgos ambientales para la salud pública en España
\cite{romanello2023}. La evidencia científica ha establecido una relación
dose-response entre temperatura y mortalidad: por encima de un umbral térmico
crítico, la mortalidad diaria aumenta de forma aproximadamente lineal con la
temperatura máxima \cite{gasparrini2017}.

En el contexto español, el verano de 2022 supuso un punto de inflexión con más
de 11.000 defunciones atribuidas al calor \cite{momo2022}. Las temporadas
2024 y 2025 han vuelto a registrar episodios de calor extremo con impacto
significativo en la mortalidad, lo que justifica su estudio sistemático. Según
el Ministerio de Sanidad, en el verano de 2025 se estimaron 3.832 muertes
atribuibles al exceso de temperatura, un 87{,}6\,\% más que en 2024
\cite{sanidad2025}.

El Sistema de Monitorización de la Mortalidad Diaria (MoMo), operativo desde
2004 y gestionado por el Instituto de Salud Carlos III (ISCIII), proporciona
datos semanales de mortalidad observada, mortalidad esperada y exceso
atribuible a temperatura para el conjunto del territorio nacional
\cite{momo2024}. Esta fuente constituye la base empírica del presente trabajo.

\textbf{Pregunta de investigación:} ¿En qué medida el calor extremo explica el
exceso de mortalidad observado en España durante los veranos de 2024 y 2025,
y qué diferencias se detectan entre ambas temporadas?

% ============================================================
\section{Metodología y resultados obtenidos}
% ============================================================

\subsection{Fuente de datos}

Se han utilizado los datos semanales del sistema MoMo del ISCIII correspondientes
al periodo comprendido entre la semana epidemiológica 20 de 2024 y la semana
35 de 2025 ($n = 68$ semanas). Las variables disponibles son: defunciones
notificadas, defunciones observadas, defunciones estimadas como línea base,
exceso de mortalidad sobre todas las causas y defunciones atribuibles a la
temperatura. La base de datos se adjunta como archivo independiente conforme
a los requisitos de la actividad.

Para el análisis se define la \textit{temporada estival} como las semanas
epidemiológicas 22--39 de cada año (aproximadamente de junio a septiembre),
periodo en el que se concentra el riesgo por calor en España.

\subsection{Modelo 1 — Análisis descriptivo de la serie temporal}

El primer modelo consiste en un análisis descriptivo de la serie temporal
completa de 68 semanas, con especial atención a las semanas estivales.
La Figura~\ref{fig:serie} muestra la evolución del exceso de mortalidad
semanal y las defunciones atribuibles al calor.

% --- FIGURA 1: Serie temporal ---
\begin{figure}[H]
\centering
\begin{tikzpicture}
\begin{axis}[
  width=\linewidth,
  height=6cm,
  xlabel={Semana epidemiológica (2024--2025)},
  ylabel={Defunciones},
  xtick={1,10,20,30,40,50,60,68},
  xticklabels={S20/24,S29/24,S39/24,S49/24,S07/25,S17/25,S27/25,S35/25},
  xticklabel style={font=\tiny, rotate=30, anchor=east},
  legend pos=north west,
  legend style={font=\small},
  grid=major,
  grid style={dashed, gray!30},
  ymin=-1000,
]
% Exceso todas causas (invertido el orden del CSV, de sem20/2024 a sem35/2025)
\addplot[azulunir, thick] coordinates {
  (1,-200)(2,-38)(3,104)(4,177)(5,-193)(6,-73)(7,83)(8,-114)
  (9,173)(10,161)(11,217)(12,678)(13,601)(14,137)(15,-234)
  (16,-422)(17,-341)(18,-152)(19,-151)(20,23)(21,-21)(22,-99)
  (23,-114)(24,61)(25,-240)(26,-570)(27,-657)(28,-685)(29,-872)
  (30,-798)(31,-207)(32,-363)(33,279)(34,1202)(35,915)(36,818)
  (37,467)(38,66)(39,309)(40,-8)(41,-70)(42,201)(43,-55)
  (44,133)(45,-13)(46,410)(47,233)(48,-275)(49,72)(50,-45)
  (51,-112)(52,158)(53,243)(54,431)(55,331)(56,223)(57,562)
  (58,1109)(59,341)(60,278)(61,72)(62,-34)(63,195)(64,905)
  (65,772)(66,272)
};
\addlegendentry{Exceso mortalidad (todas causas)}

% Atribuibles temperatura
\addplot[red!70, dashed, thick] coordinates {
  (1,3)(2,0)(3,0)(4,9)(5,12)(6,1)(7,9)(8,37)
  (9,49)(10,80)(11,310)(12,512)(13,377)(14,477)(15,71)
  (16,6)(17,3)(18,12)(19,35)(20,12)(21,10)(22,3)
  (23,0)(24,2)(25,0)(26,1)(27,6)(28,3)(29,2)
  (30,100)(31,194)(32,169)(33,285)(34,150)(35,450)(36,496)
  (37,19)(38,175)(39,129)(40,0)(41,15)(42,10)(43,12)
  (44,63)(45,11)(46,6)(47,2)(48,2)(49,1)(50,7)
  (51,1)(52,0)(53,1)(54,10)(55,9)(56,33)(57,199)
  (58,610)(59,225)(60,141)(61,135)(62,36)(63,303)(64,936)
  (65,865)(66,62)
};
\addlegendentry{Atribuibles a temperatura}

\end{axis}
\end{tikzpicture}
\caption{Evolución semanal del exceso de mortalidad y defunciones atribuibles
a la temperatura. España, semanas 20/2024--35/2025. Fuente: MoMo, ISCIII.}
\label{fig:serie}
\end{figure}

La Tabla~\ref{tab:descriptiva} resume los principales estadísticos descriptivos
para los dos veranos analizados.

% --- TABLA 1: Estadísticas descriptivas ---
\begin{table}[H]
\centering
\caption{Estadísticos descriptivos por temporada estival (semanas 22--39).}
\label{tab:descriptiva}
\begin{tabular}{lcccc}
\toprule
\textbf{Temporada} & \textbf{$n$ sem.} &
\textbf{Media exceso} & \textbf{DE exceso} &
\textbf{Total atrib. calor} \\
\midrule
Verano 2024 & 18 & 27{,}8 & 292{,}2 & 2.008 \\
Verano 2025 & 14 & 420{,}9 & 319{,}5 & 3.660 \\
\bottomrule
\end{tabular}
\end{table}

El análisis descriptivo revela diferencias notables entre ambas temporadas.
El verano de 2025 presenta una media de exceso de mortalidad de 420{,}9
defunciones semanales, frente a las 27{,}8 del verano de 2024. La desviación
estándar elevada en ambos casos ($\geq 292$) refleja la alta variabilidad
semanal, con semanas sin exceso apreciable alternadas con picos intensos.
El pico máximo se registra en la semana 27/2025 (exceso de 1.109 defunciones)
y en la semana 31/2024 (exceso de 678 defunciones).

\subsection{Modelo 2 — Regresión lineal simple}

El segundo modelo aplica una regresión lineal simple sobre las 32 semanas
estivales del periodo de estudio, tomando las defunciones atribuibles a la
temperatura ($X$) como variable independiente y el exceso de mortalidad
sobre todas las causas ($Y$) como variable dependiente.

El modelo estimado es:
\begin{equation}
  \hat{Y} = 9{,}20 + 1{,}076\,X
  \label{eq:regresion}
\end{equation}

con un coeficiente de determinación $R^2 = 0{,}57$.

La Figura~\ref{fig:regresion} muestra el diagrama de dispersión con la recta
de regresión ajustada.

% --- FIGURA 2: Diagrama de dispersión + recta ---
\begin{figure}[H]
\centering
\begin{tikzpicture}
\begin{axis}[
  width=0.85\linewidth,
  height=6.5cm,
  xlabel={Defunciones atribuibles a temperatura},
  ylabel={Exceso de mortalidad (todas causas)},
  grid=major,
  grid style={dashed, gray!30},
  legend pos=north west,
  legend style={font=\small},
]
% Datos verano 2024
\addplot[only marks, mark=o, mark size=2pt, azulunir, opacity=0.8]
coordinates {
  (3,-200)(0,-38)(0,104)(9,177)(12,-193)(1,-73)(9,83)(37,-114)
  (49,173)(80,161)(310,217)(512,678)(377,601)(477,137)(71,-234)
  (6,-422)(3,-341)(12,-152)
};
\addlegendentry{Verano 2024}

% Datos verano 2025
\addplot[only marks, mark=square*, mark size=2pt, red!70, opacity=0.8]
coordinates {
  (0,158)(1,243)(10,431)(9,331)(33,223)(199,562)(610,1109)
  (225,341)(141,278)(135,72)(36,-34)(303,195)(936,905)(865,772)
};
\addlegendentry{Verano 2025}

% Recta de regresión
\addplot[black, thick, domain=0:950, samples=2] {9.20 + 1.076*x};
\addlegendentry{Regresión: $\hat{Y}=9{,}20+1{,}076X$}

\end{axis}
\end{tikzpicture}
\caption{Diagrama de dispersión y recta de regresión entre defunciones
atribuibles a la temperatura y exceso de mortalidad. Semanas estivales
2024--2025 ($n=32$). $R^2=0{,}57$.}
\label{fig:regresion}
\end{figure}

La pendiente estimada ($\hat{\beta}_1 = 1{,}076$) indica que, por cada
defunción adicional atribuida directamente a la temperatura, el exceso total
de mortalidad aumenta en aproximadamente 1{,}08 muertes. El valor del
intercepto próximo a cero ($\hat{\beta}_0 = 9{,}20$) es coherente con la
expectativa de exceso nulo en ausencia de impacto térmico. El coeficiente
$R^2 = 0{,}57$ señala que el modelo lineal explica el 57\,\% de la varianza
del exceso de mortalidad en semanas estivales.

% ============================================================
\section{Conclusiones y líneas futuras}
% ============================================================

\subsection{Discusión comparada de los modelos}

Ambos modelos apuntan en la misma dirección: el calor extremo tiene un impacto
estadísticamente apreciable sobre la mortalidad en España, con el verano de
2025 sensiblemente más severo que el de 2024. Sin embargo, presentan
discrepancias relevantes que enriquecen el análisis.

El \textbf{modelo descriptivo} revela que varias semanas del verano de 2024
presentan excesos de mortalidad negativos (observadas $<$ esperadas), lo que
sugiere un efecto de ``recolección de mortalidad'' tras episodios previos de
exceso: las personas más vulnerables fallecen durante el pico de calor,
reduciendo las defunciones en semanas posteriores. Este fenómeno no es captado
por la regresión lineal.

La \textbf{regresión lineal} (Modelo 2), con $R^2 = 0{,}57$, confirma una
asociación positiva y sustancial, pero el 43\,\% de varianza inexplicada
apunta a que otros factores —humedad relativa, temperatura nocturna, acceso a
climatización, vulnerabilidad socioeconómica— modulan el impacto letal del
calor. Además, la relación entre temperatura y mortalidad no es estrictamente
lineal, sino que sigue una curva en J con umbral, lo que limita la validez
del modelo lineal simple para valores extremos.

\subsection{Conclusiones principales}

\begin{itemize}
  \item El verano de 2025 registró 3.660 defunciones atribuibles al calor,
        un 82\,\% más que las 2.008 del verano de 2024.
  \item La regresión lineal estima que por cada muerte directamente atribuida
        al calor el exceso total de mortalidad aumenta en 1{,}08 defunciones,
        indicando un efecto amplificador que incluye muertes asociadas pero no
        directamente codificadas como golpe de calor.
  \item Ninguno de los dos modelos por sí solo ofrece una imagen completa:
        el descriptivo aporta contexto temporal y el modelo lineal cuantifica
        la asociación, pero ambos son complementarios y necesarios.
\end{itemize}

\subsection{Limitaciones y líneas futuras}

El principal limitante del presente trabajo es la ausencia de datos de
temperatura (AEMET) en la base de datos utilizada. Incorporar la temperatura
máxima diaria permitiría construir modelos de regresión más precisos
—incluyendo términos cuadráticos o modelos de Poisson con funciones de retardo
distribuido (DLNM)— que son el estándar en epidemiología ambiental
\cite{gasparrini2017}. Como línea futura prioritaria se propone replicar
el análisis a nivel provincial para identificar territorios de mayor
vulnerabilidad, y ampliar el horizonte temporal a la serie histórica
2015--2025 para evaluar tendencias a largo plazo en el contexto del
cambio climático.

% ============================================================
% BIBLIOGRAFÍA
% ============================================================
\bibliographystyle{plain}
\begin{thebibliography}{9}

\bibitem{gasparrini2017}
Gasparrini, A., \textit{et al.} (2017).
Projections of temperature-related excess mortality under climate change
scenarios. \textit{The Lancet Planetary Health}, 1(9), e360--e367.
\url{https://doi.org/10.1016/S2542-5196(17)30156-0}

\bibitem{momo2022}
Instituto de Salud Carlos III (2022).
\textit{Informe MoMo. Exceso de mortalidad atribuible a las altas
temperaturas. Verano 2022}. ISCIII, Madrid.
\url{https://momo.isciii.es}

\bibitem{momo2024}
Instituto de Salud Carlos III (2024).
\textit{Sistema de Monitorización de la Mortalidad Diaria (MoMo).
Panel de datos}. ISCIII, Madrid.
\url{https://momo.isciii.es/panel_momo/}

\bibitem{romanello2023}
Romanello, M., \textit{et al.} (2023).
The 2023 report of the Lancet Countdown on health and climate change.
\textit{The Lancet}, 402(10419), 2022--2059.
\url{https://doi.org/10.1016/S0140-6736(23)01859-7}

\bibitem{sanidad2025}
Ministerio de Sanidad (2025).
\textit{Informe final del Plan Nacional de actuaciones preventivas de los
efectos del exceso de temperaturas sobre la salud. Verano 2025}.
Gobierno de España.
\url{https://www.sanidad.gob.es}

\end{thebibliography}

\end{document}